\chapter{Introduzione}
\label{cap:introduzione}
\intro{L'esperienza dello stage è l'elemento che porta alla conclusione del percorso accademico dello studente; pertanto è da considerare una delle prove più formanti e impegnative di tutto l'arco di studi.\\
Durante questo periodo lo studente ha l'occasione di mettere alla prova le conoscenze acquisite negli anni, ma soprattutto ha la possibilità di valutare le proprie abilità di autoapprendimento, scontrandosi con tecnologie e problemi mai visti in precedenza. 
L'introduzione al mondo del lavoro è fondamentale per rendere il più lineare possibile la transizione da studio ad applicazione pratica delle conoscenze, fornisce inoltre un assaggio delle responsabilità che gravano su un programmatore durante la permanenza nella propria azienda.
In questo capitolo viene presentata l'azienda Spazio Dev presso la quale lo studente ha svolto lo stage nel periodo che va dal giorno 23/09/2024 al 14/11/2024.}

%
%\noindent Esempio di utilizzo di un termine nel glossario \\
%\gls{api}. \\
%
%\noindent Esempio di citazione in linea \\
%\cite{site:agile-manifesto}. \\
%
%\noindent Esempio di citazione nel pie' di pagina \\
%citazione\footcite{womak:lean-thinking} \\
%
\section{Organizzazione del testo}

\begin{description}
    \item[{\hyperref[cap:processi-metodologie]{Il secondo capitolo}}] descrive i processi e le metodologie applicate dallo studente durante lo sviluppo del progetto.
    
    \item[{\hyperref[cap:descrizione-stage]{Il terzo capitolo}}] contiene tutte le informazioni relative al progetto da realizzare.
    
    \item[{\hyperref[cap:analisi-requisiti]{Il quarto capitolo}}] sviluppa la fase di analisi dei requisiti.
    
    \item[{\hyperref[cap:strumenti-tecnologie]{Il quinto capitolo}}] tratta gli strumenti e le tecnologie adottate dallo studente.
    
    \item[{\hyperref[cap:teoria]{Il sesto capitolo}}] fornisce una base teorica.
    
    \item[{\hyperref[cap:progettazione]{Il settimo capitolo}}] approfondisce l'architettura del progetto.
    
    \item[{\hyperref[cap:conclusioni]{L'ottavo capitolo}}] mostra lo stato di completamento e le future migliorie implementabili nel progetto.
\end{description}

Per quanto riguarda la stesura del documento sono state adottate le seguenti convenzioni tipografiche:
\begin{itemize}
	\item I termini che necessitano di una definizione in quanto facenti parte del gergo tecnico sono definiti all'interno del glossario.
	\item La prima occorrenza dei termini contenuti nel glossario viene evidenziata nella maniera seguente: \emph{esempio}\glsfirstoccur.
	\item I termini in lingua inglese vengono scritti in \emph{corsivo}.
\end{itemize}

\section{L'azienda}
Spazio Dev\footcite{site:spaziodev} è una \emph{software house} italiana nata a Tombolo (PD) nel 2023; nonostante la giovane età è riuscita ad affermarsi nel suo campo dando vita a molteplici progetti ambiziosi. L'obiettivo primario dell'azienda è quello di supportare diverse imprese tramite soluzioni digitali e tecnologiche sviluppate in base alle necessità.
L'azienda conta 16 membri che ricoprono i diversi ruoli di sviluppatori, grafici, responsabili marketing e sistemisiti. L'organico ha un'età media inferiore ai trenta anni, questo fattore contribuisce a rendere l'ambiente lavorativo dinamico e conviviale, perfetto per coltivare nuovi talenti.
Spazio Dev affianca e aggiorna il cliente tramite una comunicazione costante, necessaria per adattare il prodotto alle esigenze richieste. Il focus principale dell'azienda è infatti quello di fornire prodotti specifici per ciascun cliente e non la vendita di prodotti standard.\\
Il logo dell'azienda è rappresentato in Fig.~\ref{fig:logo-spaziodev}

\begin{figure}[!h] 
  \centering 
  \includegraphics[width=0.6\columnwidth]{introduzione/spazioDev-logo.png} 
  \caption{Logo dell'azienda}
  \label{fig:logo-spaziodev}
\end{figure}



\section{Servizi offerti}
L'offerta di Spazio Dev è molto ampia e consente allo staff di aiutare i clienti nei seguenti ambiti:
\begin{itemize}
  \item \textbf{Sviluppo di siti web}: La produzione di siti web è uno dei servizi più richiesti all'azienda e ne rappresenta uno dei motori trainanti. Gli sviluppatori hanno un occhio di riguardo per le nuove tecnologie, che sfruttano per offrire al cliente un prodotto accattivante e ottimizzato in maniera tale da garantire numerose visite.
  \item \textbf{Sviluppo software e integrazione IA}: I software prodotti sono solitamente indirizzati a una clientela variegata e possono spaziare in vari campi, in base alle esigenze di ciascun acquirente. Gli applicativi più venduti hanno come scopo principale la gestione dell'azienda e delle relazioni con i clienti. Per differenziarsi dalla concorrenza Spazio Dev ha introdotto nuove funzionalità impiegando modelli di IA capaci di predire \emph{trend} futuri di acquisto e  di anticipare i bisogni dei clienti.
  \item \textbf{Marketing}: Gli esperti di marketing creano contenuti di alta qualità, pertinenti e ottimizzati, atti a migliorare la presenza \emph{online} del brand. Per aumentare l'\emph{engagement} vengono studiate apposite campagne social e di posta elettronica.
\end{itemize}

\section{Clientela}
La clientela di Spazio Dev è molto variegata ed è suddivisibile nelle seguenti categorie:
\begin{itemize}
  \item Imprese di varie dimensioni specializzate in commercio e manifattura che necessitano di migliorare la propria presenza nel web, la gestione della contabilità, la logistica e le risorse umane.
  \item \emph{Startup} in cerca di visibilità e di strategie di marketing.
  \item Liberi professionisti come avvocati, medici, eccetera.
\end{itemize}



\section{Prodotto di punta}
RelAI\footcite{site:relai} (Fig.~\ref{fig:logo-RelAI}) rappresenta in tutto e per tutto l'azienda poiché nasce dall'unione di tutti i suoi \emph{core business}: sviluppo web, integrazione IA e marketing.
Il software è un \emph{CRM (Cusotmer Relationship Manager)}\glsfirstoccur  in \emph{cloud} che consente di personalizzare le relazioni con i clienti e di incrementare le vendite grazie a campagne di comunicazione su misura.
La sua architettura modulare e le numerose configurazioni si adattano perfettamente in molti settori. L'IA implementata è utile per predire futuri \emph{trend} di acquisto fondamentali per la creazione di strategie di marketing efficaci. 


\begin{figure}[!h] 
  \centering 
  \includegraphics[width=0.6\columnwidth]{introduzione/relAI-logo.png} 
  \caption{Logo dell'applicazione RelAI}
  \label{fig:logo-RelAI}
\end{figure}


